\documentclass[12pt]{article}
\usepackage{ctex}
\usepackage{amsmath, amssymb}
\usepackage{geometry}
\geometry{a4paper, margin=2.5cm}
\usepackage{enumitem}
\title{第 8 章\ 机械波\ 例题}
\author{}
\date{}

\begin{document}
\maketitle

\begin{enumerate}[label=\textbf{例\arabic*.}, leftmargin=*]

% 例 1
\item 一平面简谐波沿 $x$ 轴正方向传播，已知其波函数为
\[
  y = 0.04\cos\pi(50t-0.10x)\;\mathrm{m}.
\]
求：(1) 波的振幅、波长、周期及波速；(2) 质点振动的最大速度。

% 例 2
\item 一平面简谐波沿 $x$ 轴正方向传播，已知 $A=1.0\,\mathrm{m}$，$T=2.0\,\mathrm{s}$，$\lambda=4.0\,\mathrm{m}$。$t=0$ 时坐标原点处质点位于平衡位置沿 $y$ 轴正向运动。求：(1) 波动方程；(2) $t=1.0\,\mathrm{s}$ 时波形方程并画出波形图；(3) $x=1.0\,\mathrm{m}$ 处质点的振动方程并画出振动图。

% 例 3
\item 简谐波在 $t=0$ 时刻的波形图如右下图所示（$y$ 轴刻度为 $\frac{\sqrt{2}}{2}A$；在 $x=100\,\mathrm{m}$ 处 $y=-A$；波长 $\lambda=200\,\mathrm{m}$），设此简谐波的频率为 $250\,\mathrm{Hz}$，若波沿 $x$ 轴负方向传播，求：(1) 该波的波动方程；(2) 画出 $t=T/8$ 时刻的波形图；(3) 距原点 $O$ 为 $100\,\mathrm{m}$ 处质点的振动方程与振动速度表达式。

% 例 4
\item 频率为 $100\,\mathrm{Hz}$，传播速度为 $300\,\mathrm{m/s}$ 的平面简谐波，波线上两点振动的相位差为 $\pi/3$，则此两点相距（\quad）

% 例 5
\item 一列平面简谐波沿 $Ox$ 轴正方向传播，$t=0$ 时刻的波形图如图（$A=0.10\,\mathrm{m}$，$u=20\,\mathrm{m/s}$，$\lambda=20\,\mathrm{m}$，$P$ 点在 $x=15\,\mathrm{m}$ 处且 $y_P=0.05\,\mathrm{m}$ 处向上运动），则 $P$ 处质点的振动方程为（\quad）

% 例 6
\item 平面简谐波 $t$ 时刻的波形如图所示，此波波速为 $u$，沿 $x$ 方向传播，振幅为 $A$，频率为 $f$。$B$ 为反射点且为波节，$D$ 点位于波前零点处，且 $D$ 距 $B$ 约为 $3\lambda/4$。(1) 以 $D$ 为原点，写出波函数；(2) 以 $B$ 为 $x$ 轴坐标原点，写出入射波、反射波方程；(3) 以 $B$ 为反射点求合成波，并分析波节、波腹的位置坐标。

% 例 7
\item 一警笛发射频率为 $1500\,\mathrm{Hz}$ 的声波，并以 $22\,\mathrm{m/s}$ 的速度向某一方向运动，一人以 $6\,\mathrm{m/s}$ 的速度跟踪其后。设声速 $u=330\,\mathrm{m/s}$，求该人听到的警笛发出的声音的频率以及在警笛后方空气中声波的波长。

% 例 8
\item 一简谐平面波沿 $x$ 轴正方向传播，已知其波函数为 $y=A\cos(\omega t - bx + \varphi_0)$，试求：(1) 同一时刻波线上相位差为 $\pi/3$ 的两点之间的距离；(2) 同一波线上两点，若它们的相位差为 $\pi/2$，则该两点间的距离最近为多少？

% 例 9
\item 沿 $x$ 轴正方向传播的平面简谐波在原点 $O$ 处的振动方程为 $y_0 = A\cos(\omega t + \varphi_0)$。已知波在传播过程中无能量损失，介质中距原点 $x$ 处的能流密度大小为 $J$。求：(1) 距原点 $x$ 处质点的振动方程；(2) 距原点 $x$ 处的强度 $I$；(3) 若波在传播过程中被均匀吸收，振幅随距离按 $A(x) = A_0 e^{-\alpha x/2}$ 衰减，求 $x$ 处的强度 $I(x)$。

% 例 10
\item 两列相干平面简谐波沿同一方向传播，振幅相等，均为 $A$，频率相同，初相分别为 $\varphi_1$ 和 $\varphi_2$。求在相遇点处：(1) 合振幅最大条件与合振幅表达式；(2) 合振幅最小条件与合振幅表达式；(3) 合强度表达式。

% 例 11
\item 沿 $x$ 轴正方向传播的入射波波函数为 $y_1 = A\cos[2\pi(t/T - x/\lambda)]$，在 $x=0$ 处遇波密媒质界面发生反射，反射点为波节。求：(1) 反射波波函数；(2) 合成驻波波函数；(3) 波节与波腹的位置坐标。

% 例 12
\item 一汽笛发出频率为 $f=1000\,\mathrm{Hz}$ 的声波，汽笛以 $v_S=20\,\mathrm{m/s}$ 的速度向墙面运动，墙面以 $v_W=10\,\mathrm{m/s}$ 的速度迎着汽笛运动。声速 $u=340\,\mathrm{m/s}$。求：(1) 墙面接收到的声波频率；(2) 汽笛接收到的反射波频率。

\end{enumerate}

\end{document}
